\chapter{Discussion}
\label{ch:discussion}

Chapter~5 reported that the four categories separate, that the separation is carried almost
entirely by one of them, and that category is confounded with source dataset by construction.
This chapter asks what can be concluded from that, and is organised around the order in which
an examiner is entitled to doubt it.

\section{What the separation can and cannot support}

The measured effect, $\eta^2 = 0.1068$ at $p = 1.035\times10^{-9}$, is four times the
smallest effect the design could detect. Taken on its own terms it is unambiguous: the index
assigns systematically different values to the four groups of documents.

What it does not support is the sentence the proposal was written to test. Because each
category is drawn from a single source dataset (\S\ref{sec:confound}), a separation between
categories and a separation between corpora are the same measurement here. Fear-activating
items are ISOT-fake items. Whatever distinguishes ISOT-fake from PubMed and ISOT-true --- its
framing, but equally its vocabulary, its house style, its editorial process, its era --- is
folded into the effect and cannot be separated from it by any analysis of this data.

This is a design limitation, not an analytical one. No amount of statistical work on 400
items assembled this way can attribute the difference to framing.

\subsection{The evidence that argues against a pure source artifact}

Three observations bear on it, none decisive.

First, high outrage does not separate ($d = +0.030$, $p = 0.832$) despite coming from its own
distinct source, SemEval-2019 Task~4. If the index were simply detecting provenance, every
category drawn from a different corpus should move away from the baseline. Three of four
moving, in an ordered way, is a weaker pattern than pure source detection predicts.

Second, length is controlled. Word counts sit at $167.4$, $163.5$, $163.7$ and $162.2$ with
standard deviations near $11$. The most obvious surface confound was designed out.

Third, the ordering is not arbitrary with respect to the hypothesis. Fear-activating content
sits furthest from neutral and reward-hook content sits between. A source artifact carries no
reason to order itself along the dimension the study proposed.

Against these stands the plain fact that the confound exists and is complete. The honest
position is that the result is consistent with the hypothesis and equally consistent with a
corpus artifact, and that this study cannot choose between them.

\subsection{The experiment that would decide it}

The confound is removable by design, not by analysis. Two routes are available to a
successor:

\begin{enumerate}
\item \textbf{Within-source variation.} Draw all four categories from a single source, so
provenance is held constant while framing varies. Feasible where a source contains both
plainly-reported and emotively-framed material.
\item \textbf{Crossed design.} Ensure each source appears in more than one category, making
source and category statistically separable so their contributions can be estimated
independently.
\end{enumerate}

Either would convert the present result from a suggestive to an interpretable one. Neither
requires new instrumentation; both require rebuilding the corpus.

\section{Why high outrage does not separate}

The proposal expected outrage to separate most strongly, and it separated least. Three
readings deserve stating.

The first is that it is correct. Outrage framing may engage deliberative and affective
systems in more balanced measure than threat framing does, in which case a signed difference
between the two networks is precisely the wrong instrument for detecting it: a large
symmetric response and no response are indistinguishable in a difference.

The second is a property of the source. SemEval-2019 Task~4 labels hyperpartisan publishing
at the article level. Hyperpartisan is not synonymous with outrage-framed, and a corpus
selected on partisanship may contain a substantial fraction of items that are simply
argumentative rather than emotionally charged.

The third is the encoder. TRIBE was trained on naturalistic film and podcast viewing, and is
being asked to predict responses to synthesised speech from news text. There is no reason to
expect its sensitivity to be uniform across framing styles that were absent from training.

These are not ranked, because the data does not rank them.

\section{The instrument's standing}
\label{sec:standing}

Every value in this thesis is a prediction from a released encoder. No participant was
scanned, no brain was measured, and the observable is a cortical proxy computed from
predicted activation over two ROI unions.

That distinction is load-bearing here in a way it was not when the result looked like a null.
A null in predicted activation is a weak statement whichever way the encoder behaves. A
\emph{detected separation} in predicted activation is only a statement about media if the
encoder tracks cortex at all. The result reported in Chapter~5 therefore raises the stakes on
a question this study cannot answer from its own data.

\subsection{The checkpoint problem}

The released checkpoint is loaded in its average-subject configuration, which
\texttt{tribev2/demo\_utils.py} forces at load time. A commercial laboratory has published an
audit of that configuration reporting vertex-level $r = -0.0145$ against a measured
inter-subject ceiling of $+0.0508$: anti-correlated with real cortex, across four subjects
and all seven Yeo networks.

That audit is self-published by a company selling the per-subject alternative, uses four
subjects and one stimulus, and reports magnitudes near $0.015$ in either direction. It is a
claim about sign, not about effect size, and it has not been independently replicated. It
should not be accepted uncritically, and it cannot be dismissed either.

Work toward replicating it was begun within this project. Using the publicly released
Algonauts~2025 responses, the leave-one-subject-out noise ceiling was measured directly from
four subjects over all 1000 Schaefer parcels of episode \texttt{s01e01a}, 592 timepoints:
\begin{equation}
r_{\text{ceiling}} = 0.1517, \qquad 95\%\ \text{CI}\ [0.1484,\ 0.1551].
\end{equation}
An earlier figure of $0.2247$ is withdrawn. It was produced by a notebook cell that selected
the first key in the response file, which is a different episode from the one under
prediction, and whose orientation test could not fire on an array with fewer timepoints than
parcels, so the correlation ran across the parcel axis and the reported count of 482
``parcels'' was in fact the number of timepoints. The measurement now lives in a script that
takes the episode as an argument and asserts orientation against the atlas size.
This is the quantity any encoder must be judged against, and it is measured here rather than
quoted. It is larger than the audit's $+0.0508$ by construction and not by disagreement:
averaging vertices within a parcel cancels independent noise and raises correlations, so a
parcel-level ceiling and a vertex-level ceiling are not comparable numbers. The comparison
against the checkpoint's own predictions was not completed within this project's timeline and
is stated as unfinished rather than estimated.

\subsection{What this means for building on foundation models}

A technical report describing an architecture is not a warranty about the weights that are
released under its name. The checkpoint used here is cortical-only, which the proposal did
not anticipate; it is loaded in a configuration the accompanying paper does not foreground;
and its agreement with real cortex is contested in the literature. None of that is visible
from the paper.

The practical lesson generalises beyond this project: anyone building a measurement on a
released model should verify the checkpoint's properties from the artifact itself, as
\texttt{scripts/verify\_tribe\_checkpoint.py} does here, rather than inheriting them from the
publication.

\section{The coupling, and what the bound does instead}

The coupling between the index and the physics model's field is unidentified
(\S\ref{sec:rq3}): significant in sample, negative in held-out $R^2$, and inversely dependent
on a $\beta J$ that nothing in this study measures. Reporting a value would report a choice.

The bound is what survives. Because it requires only the observable's measured spread and the
model's own spinodal, $\alpha \ge h_c(\beta J)/\Delta X$ converts $\Delta X = 0.1241$ into a
statement no fit is needed to make: a coupling below $4.2938$ at $\beta J = 2$ cannot drive
the transition regardless of how the calibration is performed. This inverts the usual
practice in the sociophysics literature, where the field is assigned and the population is
shown to tip, without establishing that any real influence reaches the assigned magnitude.

That inversion is the contribution most likely to be reusable by others, because it applies
to any proposed field observable and can be evaluated before a corpus is collected.

\section{Session reliability, and what it permits}
\label{sec:reliability}

A second GPU session scanned 375 of the 400 items with identical text, identical code and
identical ROI definitions. That pair is a test-retest measurement, and it is reported here
because every effect in Chapter~5 otherwise rests on a single session with no statement of
how much of an item's score is the item and how much is the session.

Agreement between the two sessions, on the 375 items they share:

\begin{table}[htbp]
\centering
\begin{tabular}{lrrrr}
\hline
& $r$ & ICC(2,1) & $\mathrm{sd}(\Delta)$ & $\mathrm{sd}(\Delta)/\mathrm{sd}_{\text{between}}$ \\
\hline
$\mathrm{NAA}_{\text{signed}}$ & $0.8812$ & $0.8757$ & $0.01006$ & $0.475$ \\
$A_{\text{aff}}$ & $0.8768$ & $0.8696$ & $0.01210$ & $0.487$ \\
$A_{\text{del}}$ & $0.8575$ & $0.8544$ & $0.01285$ & $0.520$ \\
\hline
\end{tabular}
\caption{Agreement between the two scanning sessions over the 375 items both cover. ICC is reported beside $r$ because a session that shifted every score by a constant would still correlate perfectly.}
\label{tab:test-retest}
\end{table}

The intraclass correlation is reported alongside $r$ because a session that shifted every
score by a constant would still correlate perfectly, and a shift is precisely what a second
session can produce. The two agree closely here, so the disagreement is scatter rather than
offset.

\textbf{The separation replicates.} Recomputing the four-category ANOVA on the second session
over the same 375 items gives $\eta^2 = 0.0839$, $F = 11.328$, $p = 3.988\times10^{-7}$,
against $\eta^2 = 0.1054$, $F = 14.576$, $p = 5.397\times10^{-9}$ for the first session
restricted to those items. Both clear the smallest effect this design could detect,
$\eta^2 = 0.0268$. The separation reported in Chapter~5 is therefore not an artifact of one
session. The second session's effect is roughly a fifth smaller, which is the attenuation
that additional measurement error produces in a ratio of between-group to total variance.

\textbf{Per-item direction is not reliable.} The signed index is read as a direction, and
$49$ of $375$ items, $13.1\%$, reverse that direction between sessions. A statement that a
particular item is affective-leading is therefore not supportable by this instrument at this
precision, and none is made: no per-item verdict appears in the analysis, in the figures or
on the public interface. Group means over 100 items average this scatter down by an order of
magnitude, which is why the category-level result survives while the item-level one does not.

Both sessions are reported as measured. Neither is corrected toward the other and the two are
not averaged, because an average of two sessions would carry a precision that neither run
has.

\textbf{Scan order does not carry the separation.} The corpus took more than one session to
scan, which raises the question of whether position in the scan, and therefore session, is
doing the work attributed to category. It is not, and the reason is that items were scanned
interleaved rather than in category blocks: across thirds of the scan order the four
categories are balanced ($\chi^2 = 3.139$, $p = 0.7913$). There is a real drift, with mean
$\mathrm{NAA}_{\text{signed}}$ differing across thirds at $\eta^2 = 0.0167$, $p = 0.0353$.
Because it is orthogonal to category by construction, it adds scatter rather than bias:
removing the scan-third means before the category ANOVA moves the effect from
$\eta^2 = 0.1068$ to $\eta^2 = 0.1048$, $p = 1.588\times10^{-9}$. Had the corpus been scanned
one category at a time, this drift would have been indistinguishable from the finding.

\section{Limitations}

Stated here in full rather than distributed through the text.

\begin{enumerate}
\item \textbf{Category is confounded with source dataset.} The central limitation. Discussed
in \S\ref{sec:confound} and above.

\item \textbf{Single-session scores are noisy.} Test-retest ICC is $0.8757$ and the
session-to-session standard deviation is $47.5\%$ of the between-item standard deviation
(\S\ref{sec:reliability}). Category-level conclusions survive this; per-item ones do not.

\item \textbf{Predicted, not measured, activation.} No brain was scanned. Whether the
encoder tracks cortex is unresolved and contested.

\item \textbf{Cortical proxy, not a subcortical measurement.} The checkpoint is
cortical-only. The structures the original proposal named cannot be measured with it, and no
claim about them is made anywhere in this document.

\item \textbf{Average-subject configuration.} Forced by the loader. Whatever individual
variation the model can express is unavailable.

\item \textbf{Out-of-distribution stimuli.} The encoder was trained on naturalistic
audiovisual material and is fed synthesised speech from written text.

\item \textbf{The ratio form fails on $17.25\%$ of items}, and the reformulation to a signed
difference, while necessary, changes what the index means: it measures which network leads,
not by what proportion.

\item \textbf{English-language corpus from five specific sources, none of them Kenyan or
African.} A search for openly licensed Kenyan or African news data carrying usable labels
returned nothing suitable, so every item in the corpus originates from American or European
publishing. The measured spread is a property of that corpus, not of media in general, and
the bound inherits the limitation. Whether the instrument behaves the same way on media a
Kenyan reader actually encounters is untested, and assembling a labelled local corpus is the
most obvious extension of this work.

\item \textbf{Mean-field physics.} All-to-all coupling, one homogeneous population, a static
field, binary opinions. Network structure and heterogeneity generally lower $h_c$, so the
bound is conservative within the model and is not a universal statement.
\end{enumerate}

\section{What would be required next}

In order of what each would buy.

\textbf{A crossed or within-source corpus} would make the separation interpretable. It is the
cheapest of these and the one that most changes what can be claimed.

\textbf{Completing the validation} would establish whether the instrument measures anything
about brains. The noise ceiling is measured and the evaluation code is written and tested;
what remains is running the checkpoint over a held-out stimulus and correlating. It needs no
large compute.

\textbf{A subcortical checkpoint} would let the original proposal's question be asked as
posed. That means the training grid, the four source fMRI corpora and compute well outside
the scope of an undergraduate project, and it is recorded as future work rather than a plan.
